In this section, we present \textbf{\BenchmarkName}, a re-annotated instantiation of the M$^3$-Med\cite{m3med} corpus generated via our proposed pipeline. 
% 在本节中，我们介绍 \BenchmarkName，它是通过我们提出的流程生成的 M$^3$-Med\cite{m3med} 语料库的重新标注实例。
This dataset serves not only as a standalone resource but as a proof-of-concept validation for the efficacy of our Knowledge-Graph-guided synthesis methodology. 
% 该数据集不仅可以作为独立的资源，还可以作为我们基于知识图谱的合成方法有效性的概念验证。
Adopting the taxonomy established in M$^3$-Med, we categorize queries based on their topological complexity within the knowledge graph: queries corresponding to single-hop traversals are designated as \textbf{Simple Questions}, while those necessitating multi-hop reasoning paths are classified as \textbf{Complex Questions}. 
% 我们沿用 M$^3$-Med 中建立的分类体系，根据查询在知识图谱中的拓扑复杂性对其进行分类：单跳遍历的查询被定义为简单问题，而需要多跳推理路径的查询则被归类为复杂问题。
We systematically characterize the statistical and semantic properties of each category to ensure alignment with existing benchmarks.
% 我们系统地刻画了每个类别的统计和语义属性，以确保与现有基准保持一致。

While \BenchmarkName~ is derived from a curated selection of high-fidelity medical instructional videos, our primary objective extends beyond the release of a static data artifact; rather, we aim to validate the robust generative capabilities and architectural soundness of the underlying system. 
% 虽然 \BenchmarkName 源自精选的高保真医学教学视频，但我们的主要目标并非发布静态数据；相反，我们旨在验证底层系统的稳健生成能力和架构合理性。
Specifically, we seek to empirically demonstrate that our automated pipeline synthesizes query workloads exhibiting fine-grained temporal selectivity, rigorous multi-hop logical complexity, and precise visual-semantic alignment—metrics that we show to be comparable to, or exceeding, those of manually curated expert benchmarks. 
% 具体而言，我们旨在通过实证研究证明，我们的自动化流程能够合成具有细粒度时间选择性、严格的多跳逻辑复杂性和精确的视觉语义对齐的查询工作负载——这些指标与人工精心整理的专家基准测试相当，甚至更优。
Furthermore, through a structural analysis of the generated query graphs, we highlight the pipeline's potential for horizontal scalability. 
% 此外，通过对生成的查询图进行结构分析，我们重点展示了该流程的横向扩展潜力。
This capability paves the way for generating large-scale neuro-symbolic workloads where the marginal cost of annotation approaches zero, decoupling dataset scale from human labor constraints.
% 这种能力为生成大规模神经符号工作负载铺平了道路，使得标注的边际成本接近于零，从而将数据集规模与人力投入的限制解耦。
